\scanpage{44}
\unsure{X|x}xxxxxxxxxxxxxxxxxxxx

% Long oblique cancellation across the sheet; shorter strokes through individual entries.
\deleted{est réduit et \inserted{que les corps composants de} son anneau total des fractions \deleted{est} sont des extensions
séparables de $K$. Lorsque $A$ est un anneau de valuation discrète, donc $\widehat A$
également, ces conditions signifient donc aussi que le corps des fractions $\widehat K$ de $\widehat A$ est une extension séparable de $K$, condition automatiquement
satisfaite en particulier si $K$ est de caractéristique nulle.}

% The bibliographical line is encircled, with short diagonal slashes.
\deleted{IV p.224, dans la Bibliographie 36 remplacer Amer. J. par Annals.}

\deleted{IV 8.10.3. supprimer la co\unsure{n|m}dition 4\textordmasculine) , qui est inutile.}

% Hand-drawn outline encloses both lines.
\deleted{IV 8.12.8. lire: Soient $Y$ un schéma quasi-compact et quasi-séparé,
$f\!:\!X\quad Y$ ....}

\margin{\handwritten{10.9.11.1}}
\deleted{IV, p.104, \unsure{l.|1.} -10 , au lieu de (17.9.7) lire (17.9.6).}

% The first sentence is enclosed in pencil. Blank relation positions are left open.
\deleted{IV 12.1.1. (iv) ajouter qu'on y suppose $k\quad 1$. Ajouter une condition
(vi bis) L(ensemble des $x\quad X$ tels que $\mathrm{prof}^{\cdot}_{x}(\underline F_{f(x)})\quad k$ .}

\deleted{\unsure{111111111|lllllllll} Dém. La \deleted{condition} question étant locale, on peut supposer $X$ et $Y$
affines, donc $X$ est un sous-schéma fermé d'un $Y\ t_1,\ldots,t_n$ , et on est
ramené au cas où $X=Y\ t_1,\ldots,t_n$ . Utilisant alors (10.8.3) et ( \qquad ),
on trouve que $\mathrm{prof}^{\cdot}_{x}(\underline F_{f(x)})\quad k$ équivaut à $\underline{\mathrm{Ext}}^{i}(\underline F_{f(x)},\underline O_{X_{f(x)}})_x=0_x$
pour $i\quad n-k$, et on conclut grâce à (12.3.4) (qui, comme le lecteur le
constatera sans surprise, n'utilise pas dans sa démonstration le résultat sur la profondeur rectifié qu'on se propose d'établir).}

\deleted{IV 12.1.4. \unsure{Ajouter un}}

\deleted{\uline{Corollaire} 12.1.4.1. Soient $f\!:\!X\quad Y$ et $\underline F$ comme dans 12.1.1, et soit}

\deleted{IV \unsure{14.2.13|14.4.13} p. \unsure{212|202} l.11 14.4.8. a\textquotedbl) L'équivalence de cette propriété}

\deleted{IV 14.4.8. a\textquotedbl)b)Pour formuler ce\deleted{\unsure{tt}}s hypothèses,il faut supposer que
\deleted{avec les autres}
l'ensemble des composantes irréductibles de $Y$ est localement fini.}

\deleted{IV 14.4.10. (i), p.215 l.15 remplacer \textquotedbl quelconque\textquotedbl\ par \textquotedbl dont l'ensemble
des composantes irréductibles est localement fini\textquotedbl, et l.16 remplacer
\textquotedbl type fini\textquotedbl\ par \textquotedbl présentation finie\textquotedbl.}
% archive sheet number: 43
